\documentclass{ohm_project_description}

\projecttitle{M-APR: Handling von flexiblen Materialien mithilfe von Roboterarmen}
\projectauthor{TTZ Nürnberger Land / TH Nürnberg}
\projectdate{2026}

\begin{document}

\maketitle
\thispagestyle{fancy}

Flexible Materialien wie Folien, Textilien, Kunststoffe, Kabel oder biegeschlaffe Bauteile gewinnen in zahlreichen Industriezweigen zunehmend an Bedeutung. Mit steigender Produktionsautomatisierung rückt die zuverlässige und präzise Handhabung dieser Materialien in den Fokus, da herkömmliche Robotersysteme primär für starre Werkstücke ausgelegt sind.

Dieses M-APR-Projekt (3 Semester, Start WiSe 2026/2027) befasst sich mit der Handhabung flexibler Materialien durch Roboterarme. Der genaue Schwerpunkt wird gemeinsam festgelegt und kann sich sowohl auf \textbf{Hardware-Aspekte} (spezialisierte Greiftechnologien) als auch auf \textbf{Software-/KI-Aspekte} (Simulation, lernbasierte Regelung, Deep Reinforcement Learning, Vision-Language-Action Modelle) konzentrieren. Am Ende steht die Umsetzung in einer realen Roboterzelle am TTZ Nürnberger Land.

\section*{Arbeitspakete (Auswahl nach Schwerpunkt)}
\begin{itemize}[leftmargin=0.5cm]
    \setlength\itemsep{.1em}
    \item Analyse von Greifprinzipien für flexible Materialien (Vakuum, adhäsiv, formadaptiv)
    \item Entwicklung und Fertigung eines spezialisierten Greifers (Hardware-Fokus)
    \item Aufbau einer Simulationsumgebung für deformierbare Objekte (Software-Fokus)
    \item Modellbasierte Regelung oder Deep Reinforcement Learning (DRL)
    \item Vision-Language-Action (VLA) Modelle für flexible Manipulation
    \item Sim-to-Real-Transfer und Evaluation in der Roboterzelle am TTZ
\end{itemize}

\section*{Voraussetzungen}
\begin{itemize}[leftmargin=0.5cm]
    \setlength\itemsep{.1em}
    \item Bachelorabschluss in EI, Informatik, Mechatronik, Robotik o.\,ä.
    \item Programmierkenntnisse in C++ und/oder Python
    \item Grundkenntnisse in Machine Learning und Robotik
    \item Eigenständiges und zielgerichtetes Arbeiten
\end{itemize}

\vspace{0.5cm}
Das Thema ist als \textbf{M-APR (3 Semester)} konzipiert, kann nach Abstimmung auch als Projekt- oder Masterarbeit bearbeitet werden.

\vfill
\textcolor{ohm_red}{\rule{\linewidth}{0.4mm}}
\textbf{\textcolor{ohm_red}{TTZ Nürnberger Land / Labor für mobile Robotik}} \\
\textbf{Ort:} Martin-Luther-Straße 18, 91207 Lauf an der Pegnitz \\[4pt]
\begin{tabular}{@{}ll}
\textbf{Betreuer:}   & Prof. Dr. Christian Pfitzner \\
\textbf{E-Mail:}     & \href{mailto:christian.pfitzner@th-nuernberg.de}{christian.pfitzner@th-nuernberg.de} \\
\textbf{Co-Betreuer:} & M.\,Sc. Patrick Fußy \\
\textbf{E-Mail:}     & \href{mailto:patrick.fussy@th-nuernberg.de}{patrick.fussy@th-nuernberg.de} \\
\end{tabular}

\end{document}
